% Slajd 6 – architektura i skalowanie (kryteria: potencjał wdrożeniowy i skalowanie, 20 %; RULES: Design / Technical).
\begin{frame}{Architektura: pozyskiwanie danych oddzielone od prezentacji}
  \centering
  \begin{tikzpicture}[x=1mm,y=1.32mm,
      box/.style={draw=accLineStrong,fill=accSoft,rounded corners=1.2mm,text width=23mm,align=center,
                  font=\fontsize{6.3}{7.6}\selectfont,execute at begin node={\hyphenpenalty=10000},
                  inner sep=1mm,minimum height=8.5mm},
      tall/.style={box,minimum height=42mm},
      head/.style={font=\scriptsize\bfseries,text=accNavy,anchor=south},
      ext/.style={draw=accLineStrong,dashed,fill=white,rounded corners=1.2mm,text width=142mm,align=center,
                  font=\fontsize{6.3}{7.6}\selectfont,execute at begin node={\hyphenpenalty=10000},
                  inner sep=1.2mm},
      arr/.style={-{Stealth[length=1.5mm]},accNavy,line width=0.5pt},
      both/.style={{Stealth[length=1.5mm]}-{Stealth[length=1.5mm]},accNavy,line width=0.5pt}]
    % nagłówki kolumn
    \node[head] at (12.5,46.5) {ŹRÓDŁA};
    \node[head] at (42,46.5) {POZYSKANIE I~AKTUALIZACJA};
    \node[head] at (72,46.5) {MAGAZYN};
    \node[head] at (102,46.5) {API JSON};
    \node[head] at (131.5,46.5) {PREZENTACJA};
    % kolumna 1: źródła
    \node[box] (s1) at (12.5,42) {OpenStreetMap (ODbL)};
    \node[box] (s2) at (12.5,33.5) {Otwarte dane Krakowa (ArcGIS)};
    \node[box] (s3) at (12.5,25) {Właściciele obiektów};
    \node[box] (s4) at (12.5,16.5) {Społeczność: potwierdzenia, poprawki, zgłoszenia barier};
    % kolumna 2: pozyskanie
    \node[box] (i1) at (42,42) {\texttt{build\_osm\_*.py} $\rightarrow$ snapshoty w~\texttt{data/}};
    \node[box,minimum height=8mm] (i2) at (42,33) {\texttt{layers.py}: pobieranie stronami, pamięć podręczna 24~h, odświeżanie w~tle};
    \node[box,minimum height=12mm] (i3) at (42,20.5) {\texttt{api.py}: zmiany właścicieli, poprawki i~głosy, moderacja; każda zmiana trafia do historii};
    % kolumna 3: magazyn
    \node[tall] (db) at (72,29) {\textbf{SQLite} (biblioteka standardowa)\\[0.8mm]miejsca \textperiodcentered\ atrybuty: wartość, źródło, data, potwierdzenia, zaprzeczenia \textperiodcentered\ historia \textperiodcentered\ zgłoszenia barier i~głosy \textperiodcentered\ konta i~role \textperiodcentered\ zdjęcia\\[0.8mm]migracje wersjonowane, WAL, transakcje};
    % kolumna 4: API
    \node[tall] (api) at (102,29) {\textbf{\texttt{server.py} / \texttt{api.py}}\\[0.8mm]\texttt{/api/places} (profil $\rightarrow$ dopasowanie, status)\\\texttt{/api/layers} \textperiodcentered\ \texttt{/api/reports}\\\texttt{/api/ai} \textperiodcentered\ \texttt{/api/route}\\\texttt{/api/owner} \textperiodcentered\ \texttt{/api/admin}\\[0.8mm]limity rozmiaru, gzip, role};
    % kolumna 5: prezentacja
    \node[box] (p1) at (131.5,42) {Aplikacja web (telefon i~komputer), 4~języki};
    \node[box] (p2) at (131.5,33.5) {Panel właściciela};
    \node[box] (p3) at (131.5,25) {Panel administratora};
    \node[box] (p4) at (131.5,16.5) {Klienci API (planowane Open API): miasto, mapy, rezerwacje};
    % strzałki
    \draw[arr] (s1) -- (i1);
    \draw[arr] (s2) -- (i2);
    \draw[arr] (s3.east) -- (i3.west |- s3.east);
    \draw[arr] (s4.east) -- (i3.west |- s4.east);
    \draw[arr] (i1.east) -- (db.west |- i1.east);
    \draw[arr] (i2.east) -- (db.west |- i2.east);
    \draw[arr] (i3.east) -- (db.west |- i3.east);
    \draw[both] (db) -- (api);
    \draw[both] (api.east |- p1.west) -- (p1.west);
    \draw[both] (api.east |- p2.west) -- (p2.west);
    \draw[both] (api.east |- p3.west) -- (p3.west);
    \draw[arr] (api.east |- p4.west) -- (p4.west);
    % usługi opcjonalne
    \node[ext] at (72,6.5) {\textbf{Usługi opcjonalne, każda z~działaniem zastępczym:} Claude API (rekomendacje; bez klucza -- reguły słów kluczowych) \textperiodcentered\ OpenRouteService (trasy dla wózka) \textperiodcentered\ Nominatim (wyszukiwanie z~przeglądarki, tylko po zatwierdzeniu) \textperiodcentered\ backend Rampa (zgłoszenia barier przez adapter \texttt{api.js})};
  \end{tikzpicture}

  \note[item]{Najważniejsze zdanie: import i~odświeżanie danych są osobnym krokiem (skrypty, snapshoty, cache), a~prezentacja czyta tylko z~bazy przez API – można podmienić frontend albo źródło bez ruszania reszty.}
  \note[item]{Backend Rampa (FastAPI, PostgreSQL, model obserwacja $\rightarrow$ głos $\rightarrow$ zaufanie) powstał równolegle; adapter w~\texttt{api.js} pozwala przenieść zgłoszenia barier tam bez zmian w~aplikacji. To ścieżka skalowania na większe wdrożenie.}
  \note[item]{Uczciwie: konfiguracja warstw miejskich jest dziś w~kodzie (\texttt{layers.py}), nie w~pliku miasta – to pierwszy krok przy drugim mieście.}
\end{frame}
